% SOURCE_STATEMENT: "Given a compact symplectic manifold equipped both with a compatible-complex-structure quantization and with a Lagrangian-fibration quantization, construct a canonical isomorphism between the resulting quantum Hilbert spaces."
% Preparation and final self-audit: CONTRIBUTING.md prompts C and D; source/audit artifacts retained in the local physics_statements_2026_10_11 batch.
% !TeX program = lualatex
\documentclass[11pt]{article}
\usepackage[margin=25mm]{geometry}
\usepackage{fontspec}
\setmainfont{Latin Modern Roman}
\usepackage{amsmath,amssymb,mathtools}
\usepackage{unicode-math}
\setmathfont{Latin Modern Math}
\usepackage{xurl,hyperref}
\hypersetup{colorlinks=true,urlcolor=blue}
\setcounter{secnumdepth}{0}
\setlength{\parindent}{0pt}
\setlength{\parskip}{5pt}
\setlength{\emergencystretch}{3em}
\newcommand{\sref}[2]{\hyperlink{source:#1:#2}{[S#2]}}
\newcommand{\cataloguescope}[1]{\paragraph{Recorded target.}#1}
\newcommand{\R}{\mathbb R}
\newcommand{\C}{\mathbb C}
\newcommand{\N}{\mathbb N}
\newcommand{\Z}{\mathbb Z}
\newcommand{\Q}{\mathbb Q}
\newcommand{\F}{\mathbb F}
\newcommand{\eps}{\varepsilon}
\begin{document}
\section{Canonical Equivalence of Complex and Lagrangian Quantizations}\label{statement}

\subsection{Definitions and mathematical statement}

\paragraph{Shared notation.}$\mathbb N=\{1,2,\ldots\}$; the standard number fields have their usual meanings. All additional conventions are stated locally.


\paragraph{Definitions and conventions.}
A quantizable compact symplectic manifold $(X,\omega)$ has a prequantum Hermitian line bundle with compatible connection. A compatible complex polarization gives a Hilbert space of appropriate holomorphic sections. A real polarization from a Lagrangian fibration gives states associated with its Bohr--Sommerfeld fibers, on which the restricted connection has nonzero parallel sections. An isomorphism is canonical only if it is naturally determined by the specified geometric data, rather than chosen from an arbitrary basis.

\paragraph{Target retained from the canonical source catalogue.}
\paragraph{Catalogue wording retained for traceability.}
Given a compact symplectic manifold equipped both with a compatible-complex-structure quantization and with a Lagrangian-fibration quantization, construct a canonical isomorphism between the resulting quantum Hilbert spaces.

\paragraph{Source scope and qualification.}
The prequantum data, singular-fiber treatment and any half-form corrections must match. Equality of dimensions or existence of some vector-space isomorphism is weaker than the canonical comparison requested.
\subsection{Short English statement}
Construct a natural identification between the Hilbert spaces from complex and Lagrangian polarizations of the same symplectic manifold.
\subsection{Sources}
\begin{itemize}
\item[\hypertarget{source:statement:1}{S1}] ULAM.AI and the UnsolvedMath Contributors, UnsolvedMath, catalogue entry. Catalogue curation under CC BY 4.0; definitions and English exposition added locally.
\url{https://huggingface.co/datasets/ulamai/UnsolvedMath}.
\item[\hypertarget{source:statement:2}{S2}] Subfactors and Applications (2019), Oberwolfach Reports, DOI 10.4171/owr/2019/49, pp. 3053--3140. Primary report accessed 2026-10-11; relevant PDF pages 73.
\url{https://doi.org/10.4171/owr/2019/49}.
\item[\hypertarget{source:statement:3}{S3}] Hattori--Yamashita, spectral convergence in geometric quantization (catalogue-listed context; not independently consulted in this preparation).
\url{https://arxiv.org/abs/1912.07994}.
\item[\hypertarget{source:statement:4}{S4}] Hattori--Yamashita, toric symplectic manifolds (catalogue-listed context; not independently consulted in this preparation).
\url{https://arxiv.org/abs/2002.12495}.
\end{itemize}
\label{end:statement}
\end{document}
